% =====================================================================
%  R. Nielsen, SVAR TIL HR. DR. PHIL. PASTOR ZEUTHEN I ANLEDNING AF HANS
%  SENDEBREV OM TRO OG VIDEN. Kjøbenhavn. Forlagt af den Gyldendalske
%  Boghandel (F. Hegel). Trykt hos J. H. Schultz. 1866.   16 pp.
%
%  Diplomatic transcription from the Royal Danish Library scan
%  (`bibliotek/Nielsen, Rasmus/1866-svar-til-zeuthen.pdf`,
%  https://www.kb.dk/e-mat/dod/130024834396-color.pdf, 21 PDF pages,
%  A4 page boxes, one 300 ppi colour JPEG per page,
%  sha256 9291ca32d750a36b44b1fe8a4c6727ae3201eb6209e9066110d90a7bef3c7ef4).
%
%  The exact title, from the title page (PDF 4), line by line:
%      Svar                                  [large, letterspaced]
%      til
%      Hr. Dr. phil. Pastor Zeuthen          ["Dr. phil." in antiqua]
%      i Anledning af hans Sendebrev
%      om Tro og Viden.                      [large bold Fraktur]
%      Af
%      R. Nielsen.
%      Kjøbenhavn.
%      Forlagt af den Gyldendalske Boghandel (F. Hegel).
%      Trykt hos J. H. Schultz.              [small; "Schultz" letterspaced]
%      1866.
%  KB's catalogue record: "Svar til Hr. dr. phil. Pastor Zeuthen i Anledning
%  af hans Sendebrev om Tro og Viden", Kjøbenhavn: Gyldendal, 1866, 16 s.
%  The reply is to F. L. B. Zeuthen's open letter (Sendebrev) on faith and
%  knowledge; Nielsen cites it by page, "(S. 7)", "(S. 12)".
%
%  This header is the edition's apparatus. The transcription harness
%  (pagemap.py, ocr.sh, spacing.py, check.py, splice.py, joints.py) is not
%  tracked in git and is unrunnable without the scan, so everything an
%  outside reader needs in order to check this text against the scan is
%  recorded here.
%
% ---------------------------------------------------------------------
%  1. PAGE MAP  (verified by eye 22 Sep 2026; see pagemap.py)
% ---------------------------------------------------------------------
%
%      printed 1--16   =   PDF page - 3      (PDF 4--19)
%
%  Uniform. The folio of every one of PDF 6--19 was read on the image:
%  3--16 in unbroken sequence. Two pages carry no folio: p. [1] (PDF 4,
%  the title page) and p. [2] (PDF 5, the opening of the letter). The
%  title verso is therefore the first page of text: the pamphlet is a
%  single sheet of 16 pages with no blank leaf.
%      PDF 1       KB digitisation notice -- not part of the book
%      PDF 2--3    board / KB shelfmark label and barcode
%      PDF 4       TITLE PAGE = p. [1]                   [transcribed]
%      PDF 5--19   pp. [2]--16, the letter                [transcribed]
%      PDF 20--21  blank / back board
%
% ---------------------------------------------------------------------
%  2. TYPE, STRUCTURE AND CONVENTIONS
% ---------------------------------------------------------------------
%
%  FRAKTUR throughout. Letterspacing (Sperrsatz) is the emphasis ->
%  \emph{} (38 places). Latin and other foreign words set in ANTIQUA ->
%  \textit{} ("Dr. phil.", natura, creatura, bonum est, quia deus vult,
%  testimonium spiritus). The Greek on p. 6 (κτισις, φυσις) is printed
%  WITHOUT accents or breathings, and is typed so.
%  Structure: an open letter, no sections. A centred bold salutation
%  "Høistærede Hr. Doctor!" on p. [2]; the text runs to p. 16 and ends
%  with a short centred rule -- no signature, place or date.
%  Run-in numbered points 1) 2) 3) (p. [2]) stay run-in.
%  One footnote, p. 15, marked *) in the print.
%  Orthography as printed (1866); ø throughout -- no plain o for ø was
%  found, and no ö. aa, never å. Quotation marks „…“ as printed. Fraktur
%  I/J: house convention, I for I-words, J for J-words. Printed line
%  breaks are kept; words broken at a line end are closed up with \-% .
%  NOT TRANSCRIBED: folios, sheet signatures, library stamps.
%
% ---------------------------------------------------------------------
%  3. PRINTER'S DEFECTS  (transcribed as printed, logged at the site)
% ---------------------------------------------------------------------
%
%  p. [2] Natuvidenskaben; p. 6 traditonelle; p. 8 egenlig; p. 9
%  Naturpocessen; p. 11 Dogmamatiks (syllable doubled over the line
%  break); p. 13 Grædselinie; p. 15 "o s. v." (no stop after o).
%  Quotation marks: unbalanced quotations on pp. 9, 10 and 13 cancel
%  out, so the count is 102/102 -- each is logged at its site.
%  Doubtful readings, sense-reading kept on positive evidence: title
%  "Trykt" (damaged r, reads "Tiykt"), p. 4 "naar" (ink speck), p. 15
%  "Kirkens og" (faint mark, not a hyphen).
%
% ---------------------------------------------------------------------
%  4. ACCURACY
% ---------------------------------------------------------------------
%
%  Word-accuracy: every batch diffed against the scan's embedded KB/ABBYY
%  text layer before splicing (ocrdiff.py). 1 batch, 37 candidates, 0
%  transcription errors corrected, 0 unresolved (all 37 the layer's ø/æ
%  losses, one out-of-order phrase, one repeat). All 16 pages rendered and
%  read at the image by the transcribing agent. Caller cross-check
%  22 Sep 2026: a page-by-page word alignment of the whole text against
%  BOTH witnesses (the embedded layer and tesseract Fraktur) flagged one
%  place where the two agreed against the transcription (p. 10 "et
%  uegentligt Mirakel", both witnesses "er"); settled at the image as
%  "et" -- the transcription is right. One emphasis call (p. 10
%  "Eiendommelighed") spot-checked at the image: correct.
%  See TRANSCRIPTION-PLAYBOOK.md §3 step 4.
%
%  STATUS: COMPLETE. pp. 1--16 transcribed 2026-09-22 in one batch.
% =====================================================================
\documentclass[12pt,a4paper]{article}
\usepackage[utf8]{inputenc}
\usepackage[T1]{fontenc}
\usepackage[danish]{babel}
\usepackage{textalpha}
\usepackage[protrusion=true,expansion=true]{microtype}
\usepackage[margin=1.3in]{geometry}
\usepackage{setspace}
\usepackage{libertinus}
\usepackage{libertinust1math}
\usepackage{graphicx}
\DeclareUnicodeCharacter{0254}{\reflectbox{c}}
\usepackage{fancyhdr}
\usepackage[colorlinks=true,linkcolor=black,urlcolor=black,
  pdftitle={Svar til Hr. Dr. phil. Pastor Zeuthen},
  pdfauthor={R. Nielsen}]{hyperref}
\setlength{\headheight}{15pt}
\pagestyle{fancy}
\fancyhf{}
\fancyhead[L]{\textit{Svar til Hr. Dr. phil. Pastor Zeuthen}}
\fancyhead[R]{\thepage}
\setlength{\parindent}{1.5em}
\setstretch{1.3}
\renewcommand{\thefootnote}{\ifcase\value{footnote}\or *)\or **)\or ***)\else
  \arabic{footnote})\fi}
\usepackage{marginnote}
\newcommand{\opage}[1]{\marginnote{\footnotesize\textit{[#1]}}}
\newcommand{\apage}[1]{\marginnote{\footnotesize\textit{[#1]}}}
\begin{document}

% --- p. 1 ---
\opage{1}\setcounter{footnote}{0}%
% Title page (PDF 4), set centred line by line. "Svar" in large letterspaced
% Fraktur; "til" small; "Dr. phil." in antiqua, the rest of that line in
% Fraktur; "om Tro og Viden." in large bold Fraktur; "R. Nielsen." bold
% Fraktur; "Kjøbenhavn." bold Fraktur; "Trykt hos J. H. Schultz." small,
% "Schultz" letterspaced; "1866." in antiqua figures. KB library stamp and
% ex-libris stamp (SG-JD-FH monogram) omitted.
% "Trykt": the r is damaged in this copy and reads like "Tiykt"; no other
% letter fits, so the sense-reading is kept (PLAYBOOK §2).
\begin{center}
{\Huge S v a r}\\[1.5em]
{\small til}\\[1.5em]
{\LARGE Hr. \textit{Dr. phil.} Pastor Zeuthen}\\[1em]
{\large i Anledning af hans Sendebrev}\\[1em]
{\Huge\bfseries om Tro og Viden.}\\[1em]
{\small Af}\\[0.5em]
{\large\bfseries R. Nielsen.}\\[6em]
{\bfseries Kjøbenhavn.}\\[0.5em]
Forlagt af den Gyldendalske Boghandel (F. Hegel).\\[0.3em]
{\footnotesize Trykt hos J. H. \emph{Schultz}.}\\[0.5em]
1866.
\end{center}
\clearpage
% --- p. 2 ---
\opage{2}\setcounter{footnote}{0}%
% p. [2] (PDF 5): no folio. The salutation is centred on its own line in
% bold Fraktur; the letter opens with a two-line initial V.
\begin{center}
{\large\bfseries Høistærede Hr. Doctor!}
\end{center}

Vi kunne, saavidt jeg skjønner, blive enige om meget
væsentlige Puncter, dersom vi blot kunde blive enige om,
at en Mirakelvidenskab, en Videnskab, der forklarer Mirakler,
hører blandt de umulige Ting. De spørger mig i Deres Sendebrev
(S. 7), hvorledes man kan udtale det „som en afgjort Sag,
at Videnskaben negter Underet“, og tilføier strax efter: „Hvad
Naturvidenskaben angaaer, da maa man vel tilstaae, at dens
Naturbegreb ikke er theologisk, og at den ingen Plads har
for det, vi i religiøs Forstand kalde det Underfulde, men er
Naturvidenskaben al Videnskab, den høieste Videnskab?“ Jeg
svarer: at der vel foruden Natuvidenskaben gives en Idee\-% "Natuvidenskaben": sic, r dropped (printer's error; both witnesses and image agree)
videnskab; men om denne sidste gjælder, efter min Syns\-%
maade, tre Ting: 1) den maa være bygget paa Naturlærens
mathematisk-empiriske Grundlag, 2) den maa, ligesaa vel som
Naturlæren være en Erkjendelse i Kraft af Love; den kan
3) ikke tænkes at staae i nogen væsentlig eller principiel Strid
med Naturlæren. Maa nu Naturlæren som Videnskab, vel at
mærke, som Videnskab, benegte Muligheden af, at „en virkelig
Død“ virkelig kan staae op af Graven (saa at der altsaa ikke
kan gives nogensomhelst naturvidenskabelig Redegjørelse for
Opstandelsens Mirakel), da maa Ideelæren, saasandt den vil
gjælde for Videnskab, negte det Samme. Naturlæren
vilde jo være falsk, dersom den ved sin Udtydning af Kjends\-%
gjerningerne erklærede det for umuligt, som dog ifølge
videnskabelig Opfattelse af Naturens Idee var muligt; omvendt
% --- p. 3 ---
\opage{3}\setcounter{footnote}{0}%
vilde den Ideelære kun være huul Speculation, som ved sin
Udvidelse af Mulighederne kom i Strid med Naturlærens
første ufravigelige Grundsætninger. At ville tilveiebringe
Forsoning mellem Religion og Videnskab ved at sætte Splid
mellem Videnskab og Videnskab gaaer paa ingen Maade an.

Den gamle Drøm, thi for Andet end en Drøm kan jeg
nu ikke ansee det, om en høieste Videnskab, der skulde være
istand til at udfinde Miraklets „Tænkelighed“ og „Mulighed“,
grunder sig, om jeg ikke feiler, paa et Par Tvetydigheder,
Tvetydigheden i Ordet \emph{Under} og Tvetydigheden i den Brug,
Theologien gjør af det Teleologiske.

Hvad er vel, Underet, naar Theologien skal til at
tractere det „objectivt“, Andet end en stor Tvetydighed?
„Alt“ --- siger en skarpsindig Theolog --- „er et Under.
Enhver Begivenhed er Guddommens Virkning, umiddelbart
grundet i den guddommelige Kraft; men den viser sig som
Naturvirkning middelbart grundet i Gud. I objectiv Forstand
er altsaa Alt Under eller Intet Under; Frøets Opvaagnen
til Vaarens Liv er et Under saavelsom den Dødes Opvaagnen
af Graven, Helbredelse af Lægemidler saavelsom ved Haands\-%
paalæggelse, og ethvert Støttepunkt for Begrebet falder hen.
Naturens System af Kræfter er et i Evighed uudgrundeligt
Dyb, og vi beskæmme os kun, om vi formaste os til at drage
Grændselinier, hvor der ingen gives“. Men at en Tvetydig\-%
hedslære, der sætter „den Dødes Opvaagnen af Graven“ og
„Frøets Opvaagnen til Vaarens Liv“ paa lige Trin i Mulig\-%
hedernes Række, ikke just er skikket til at forsone Tro og
Viden, men allerhelst maatte anvises „en egen Plads udenfor
Videnskaben“, ville vel i vor Tid de Fleste indrømme.

Til den tvetydige Brug af Ordet Under svarer i Theo\-%
logien en ligesaa tvetydig Brug af Ordene \emph{Teleologie} og
det \emph{Teleologiske}. Der er nemlig en Forskjel, ja en meget
væsentlig Forskjel imellem en Teleologi, der er det høiere
Udtryk for Lovenes Opfyldelse, en Fremstilling af Naturens
virkelige Formaal, og en phantastisk, Videnskabens Love
oversvævende Teleologie. Da det nu, efter min Mening,
hører til den philosophiske Dannelses første Fordringer, at
% --- p. 4 ---
\opage{4}\setcounter{footnote}{0}%
være orienteret angaaende Forskjellen imellem det Lovbestemte
og det Phantastiske, har jeg forlængst afhandlet de herhen
hørende Tvetydigheder i Propædeutiken, og hensætter her,
til Exempel, en lille Prøve: „A. Antager Du, at Natur\-%
forskerens Begreb om Naturloven staaer i et væsentligt eller
i et blot tilfældigt Forhold til Naturvidenskaben? C. I et
væsentligt, naturligviis! A. Du citerede ovenfor et Stykke
af H. C. Ørsteds \emph{Aanden i Naturen}, og vil da ogsaa% "Ørsteds": the Fraktur capital Ø (inner diagonal; cf. the plain O of "Ord", p. 11); "C." -- the stop is small and set low
mindes, at der i samme Skrift forekommer en Afhandling
„Naturvidenskabens Forhold til adskillige vigtige Religions\-%
gjenstande“, hvoraf første Deel, der har til Overskrift: „Natur\-%
lovenes Uforanderlighed“, begynder saaledes: „I min Bog,
Aanden i Naturen, har jeg med fuld Overbeviisning hyldet
den gamle Lære om Naturlovenes Uforanderlighed....“ Men
naar Naturlovene ere uforanderlige, hvorledes skulle vi da
udtyde Sætninger som disse: „Frøets Opvaagnen til Vaarens
Liv er et Under, saavelsom den Dødes Opvaagnen af Graven,
Helbredelsen af Lægemidler saavelsom ved Haandspaalæggelse“?
C. Den hermeneutiske Nøgle ligger jo i de Ord: „Enhver
Begivenhed er Guddommens Virken, umiddelbart grundet i
den guddommelige Kraft, men den viser sig som Naturvirkning
middelbart grundet i Gud. I objectiv Betydning er altsaa
Alt Under eller Intet Under.“ A. Hvis de Døde altsaa
vaagnede i Graven hvert Foraar, naar Frøet „opvaagner til% "naar": a speck below the n makes it look like "yaar"; the sort is n
Liv“, saa vilde det Første være i ligesaa god Overeensstem\-%
melse med „Naturlovenes Uforanderlighed“, som det Sidste?
C. Ja, hvorfor ikke, dersom en saadan Opvaagnen virkelig
var \emph{teleologisk}? Naturen er jo“ --- her laanes en Sætning
fra \textit{Dr.} H. Martensens „Christelige Dogmatik“ --- „et System,
der befinder sig i en fortsat teleologisk Udvikling, en fortsat
Skabelse!“ Hvorfor skulde de Døde ikke have Lov til at
vaagne hvert Foraar, naar de kun vaagnede teleologisk, ifølge
en fortsat Skabelse? A. At Svovlarsenik kan blive til
Vand, Vand til Viin, Viin til Kul, og Kul til Guld, er
altsaa i fuldkommen Overeensstemmelse med Naturlovenes
Uforanderlighed? C. Dersom det, vel at mærke, skeer
\emph{teleologisk}, ifølge en fortsat Skabelse! A. At et Men\-%
% --- p. 5 ---
\opage{5}\setcounter{footnote}{0}%
neske vandrer paa Havet uden at synke er da ligesaa stem\-%
mende med Naturlovenes Uforanderlighed, som at En falder
i Vandet og drukner? C. Mennesket, der ellers er under\-%
kastet Tyngdens Love, kan meget vel vandre paa Havet, naar
det skeer \emph{teleologisk}, ifølge en fortsat Skabelse! A. Men
saa har det, hvad Naturlovenes Uforanderlighed angaaer, vel
heller ikke, naar vi ret besinde os, noget for den dybere
Tænkning anstødeligt, at en Pave, Leo XII, har kunnet
beatificere Minoriten Julianus, som i vort vantro, revolu\-%
tionære nittende Aarhundrede fik stegte Fugle til at flyve?
C. At jeg for mit Vedkommende antager Underet med de stegte
Fugle for en dum Fabel, en Præsteløgn, hidrører slet ikke
derfra, at jeg i Begivenheden selv seer noget Brud paa Natur\-%
lovenes Uforanderlighed, men simpelthen derfra, at jeg ikke er
Catholik og følgelig ikke kan overbevise mig om, at Begiven\-%
heden er skeet \emph{teleologisk}, ifølge en \emph{fortsat Skabelse}.
Var jeg Catholik og havde en Catholiks Tro, da skulde jeg
vel være istand til at forsone min Tro med min Viden, mit
Skabelsesbegreb med mit Naturbegreb; da skulde jeg vel i
Betragtning, ja i dogmatisk Betragtning af Teleologien i de
stegte Fugle, kunne forklare Miraklet saaledes: „Naturen“ ---
her laanes igjen en Sætning af \textit{Dr.} H. Martensens „Christe\-%
lige Dogmatik“ --- „fornegter ikke Skabelsesbegrebet, og
Underet er netop det Punkt, hvor Naturens Afhængighed af
den skabende Frihed bliver fuldkommen aabenbar.“ Vilde den
evige Fornuft, den treenige Gud, den absolute Idee, give os
et Tegn, der paa en ret iøinefaldende Maade kunde vise disse
Tiders Vantro, at Naturen, et System, der befinder sig i en
fortsat \emph{teleologisk} Udvikling, ikke fornegter Skabelses\-%
begrebet: hvor kunde den evige efter \emph{uforanderlige
Naturlove} virkende Fornuft, da vælge Tegnet mere \emph{teleo\-%
logisk} end netop dette Tegn med Minoritens stegte Fugle;
thi dette staaer dogmatisk fast: ere Fuglene stegte \emph{teleologisk},
ifølge en i Skabelsen fortsat Naturproces, saa maae de ogsaa
kunne flyve \emph{teleologisk}, ifølge en i Naturen fortsat Skabelses\-%
proces. A. Du er en ægte Theolog, Du kan forsone Tro
% --- p. 6 ---
\opage{6}\setcounter{footnote}{0}%
og Viden.“ (Forelæsninger over „Phil. Propæd.“ 1861--62;
S. 158--59).

De vil i denne korte Forhandling, Hr. Doctor, vistnok
hverken finde Fordreining eller Overdrivelse. Det er jo en
bekjendt Sag, at den catholske, til Kirkens Myndighed støttede,
paa sit middelalderlig traditonelle Grundlag fast stillede% "traditonelle": sic, i dropped (printer's error)
Scholastiker har en ganske anden Sikkerhed i at forklare og
forsvare sine Helgenmirakler end den protestantiske, af Fornuft\-%
oplysning blendede, i sytten halvvidenskabelige Hensyn hildede
Theolog har i at forklare enten Christi eller hans Apostles
Gjerninger.

For end yderligere at overbevise os om, hvad det i
Virkeligheden bliver til, naar den formeentlig høiere Viden\-%
skab giver sig ifærd med at forklare og begrunde Aabenbaringens
Undere, behøve vi blot at kaste et Blik paa „den christelige
Dogmatik“; vi have jo lykkeligviis i vor nyere Litteratur
den Dogmatik, der med Grund kan siges at repræsentere et
Høidepunkt af theologisk Videnskabelighed.

Ja, skulde nogen Videnskab være istand til at for\-%
klare Underet, saa maatte det rigtignok være Dogmatiken.
Hvad har da vor „christelige Dogmatik“ bragt frem til viden\-%
skabelig Forstaaelse af Aabenbaringens tre Grundundere:
\emph{Skabelsen, Incarnationen} og \emph{Inspirationen}?

Om \emph{Skabelsen} læres der, at Gud „frembringer en
Skabning, der tillige frembringer og udvikler sig selv i given
Selvstændighed“. Skal det være muligt at forbinde nogen
Mening med en saa dunkel Sætning, maa Meningen aaben\-%
bart være den, at Skabervirksomheden tillige er Naturvirksomhed,
at Underet tillige skeer paa naturlig Maade. Hermed stemmer
det saa rigtignok, at Verden er skabt paa den ene Side med
Nødvendighed paa den anden med Frihed, paa den ene Side
af Intet, paa den anden af Noget; at Skabelsen i Grunden
er Kosmogonie og Kosmogonien Skabelse. „Jødedommen“,
hedder det, „opfatter Verden overveiende som \textit{creatura}, ikke
som \textit{natura}, som κτισις, ikke som φυσις. Men just% Greek printed without accents or breathings (checked at 600 dpi); kappa in the curly form, typed κ
derfor er Skabelsens hele Betydning ikke opladt for Jøde\-%
dommen“; o. s. v. Er Miraklet først paa denne Maade blevet
% --- p. 7 ---
\opage{7}\setcounter{footnote}{0}%
Naturbegivenhed og Naturbegivenheden Mirakel, er Natur og
Skabelse slaaet sammen saaledes, at Naturen til en vis Grad
er skabt og det Skabte Natur, saa bliver det en let Sag at
indføre Gradationer, ikke blot i det Skabtes „givne Selv\-%
stændighed“, men i selve Skabelsesmiraklet; det ene Skabelses\-%
mirakel bliver da større, stærkere og mærkværdigere end det
andet: vi faae en dogmatisk Gradestok for Mirakler. Saa\-%
ledes seer „den christelige Dogmatik“ i Christo paa een Gang
„den menneskelige Naturs fuldendte Skabelse og den menneske\-%
blevne Gud selv“, seer med andre Ord, hvorledes den Uskabte
selv er skabt: dette er nu det allerstørste Mirakel. Skabelsen
af „den første Adam“ er jo vistnok ogsaa et stort Mirakel.
„Den første Adam er \emph{skabt} i en Betydning i hvilken ingen
af hans Efterkommere er det“; men dette store Skabelses\-%
mirakel staaer dog mange Grader under Incarnationsmiraklet.
Ligesom nu Miraklet med Skabelsen af den første Adam staaer
mange speculative Grader under Miraklet med den anden,
saaledes staaer Miraklet med Evas Skabelse, skjøndt i
sig selv et respectabelt Mirakel, idetmindste to og en halv
Grad under Miraklet med Adams Skabelse. Saamange
Grader i Miraklet Skabelse, saamange Betydninger af Ordet
Skabelse! Er den anden Adam skabt i en Betydning, i
hvilken ingen anden Skabning er det; er den første Adam
skabt i en Grad og Betydning, i hvilken ingen af hans
Efterkommere er det; er Eva skabt i en Grad og Betydning,
i hvilken ingen anden Qvinde er det; Heroerne og Genierne
igjen skabte i andre Grader og andre Betydninger end ordi\-%
naire Mennesker: saa indsees, at Skabelsens Mirakelgrader
maae kunne variere indenfor store Knudepunkter, næsten ligesaa
continuerlig som de naturlige Varmegrader. Denne flydende,
om end undertiden tyktflydende, Eenhed af Natur og Skabelse,
bliver saa tilsidst ved Tilsætning af „Traducianisme“ og
„Creatianisme“ i mærkværdig Grad tyndtflydende. „Det er
Traducianismens Sandhed, at ethvert menneskeligt Individuum
er en Affødning af Slægtens Naturvirksomhed“. Men det
er „Creatianismens Sandhed, at denne hemmelighedsfulde
Naturvirksomhed er Organ og Middel for Guds individuali\-%
% --- p. 8 ---
\opage{8}\setcounter{footnote}{0}%
serende Skabervirksomhed“. Skabelsen er hørt op, thi den er
afløst af Opholdelsen; men Opholdelsen er en fortsat Skabelse,
Skabelsen er altsaa ikke hørt op; Skabelsen er til en vis Grad
hørt op, til en vis Grad ikke hørt op. Dette er i en
Hovedsum den Forklaring, „den christelige Dogmatik“ giver os
af Skabelsens Under.

Og hvad lærer „den christelige Dogmatik“ om Incarna\-%
tionens Under? Blandt Andet Følgende: „Naar man har
fundet det selvmodsigende, at den evige Logos bliver Menneske,
„fordi den Evige og Allestedsnærværende ikke kan lade sig
føde midt i Tiden“; da kan det vistnok ikke være Opgaven at
tænke, at den evige Logos med Incarnationen skulde ophøre
at existere i sin almindelige Verdensaabenbaring, eller at Logos
som selvbevidst personligt Væsen skulde være indesluttet i
Modersliv, fødes som Barn, tage til i Kundskab o. s. v., thi
denne Forestilling ophæver selve Fødselens Begreb. Fødsel
i Tid medfører nødvendigt Forestillingen om en Fremgang
fra det Ubevidste til det Bevidste, fra Mulighed til Virkelighed,
fra Sædekorn og Spire til fuldbaaren Organisation. At den
guddommelige Logos lader sig føde, vil derfor sige, at han
indsænker sig i Menneskehedens Skjød som Mulighed, som
hellig Sæd, for at kunne opstige i Menneskeslægtens Midte i
midlende og frelsende \emph{Menneske}aabenbarelse“. Her have
vi altsaa Forklaringen. Men, Hr. Doctor, skulde vi ikke let
kunne blive enige om, at denne Forklaring selv trænger
høilig til en Forklaring? Thi hvad er det den guddom\-%
melige Logos egenlig har „indsænket“ i „Menneskehedens% "egenlig": sic (for "egentlig"; cf. "egentlige" below)
Skjød“, \emph{sig selv} eller \emph{ikke sig selv}? Er det ikke sig selv,
ikke sit egentlige Selv, saa er det jo heller ikke Logos selv,
der er incarneret. Det Væsen, hvis egentlige Selv ikke kan
„indsænkes“ som Mulighed i Modersliv, kan ikke incarneres.
Men er det virkelig sig selv, sit mulige, væsentlige Selv, men
dog ikke det fra Evighed af selvbevidste personlige Selv, hvorom
her tales: saa maa den guddommelige Logos jo have et
dobbelt Selv, et, der kan incarneres, og et, der ikke kan in\-%
carneres. Den Logos, der kan incarneres, er altsaa i sig selv,
i sit egentlige Selv, en Anden end den evige personlige Logos
% --- p. 9 ---
\opage{9}\setcounter{footnote}{0}%
selv; med andre Ord: den evige personlige Logos selv er
ifølge „den christelige Dogmatiks“ Forklaring slet ikke bleven
incarneret.

Og hvilken Udtydning giver „den christelige Dogmatik“
os saa tilsidst af Inspirationens Under? Det er Begiven\-%
heden paa Pintsefesten (Apost. Gjerng. 2 Cap. V. 1--12).
Lyden fra Himlen, det fremfarende vældige Veir, Ildtungerne,
Sproggaverne, at „de alle bleve fyldte af den Hellig Aand
og begyndte at tale med andre Tungemaal, eftersom Aanden
gav dem at tale“, der skulde forklares. „Den nye Skabelse“,
siger Dogmatiken, „staaer, hvad det Mirakuløse angaaer, i
ingen Henseende tilbage for den gamle, her gjøres „en over\-%
naturlig Begyndelse“. Men, see vi nærmere til, saa er Mi\-%
raklet ved dogmatisk videnskabelig Grademaaling --- thi Aands\-%
miraklet har naturligviis ogsaa sine Grader og Gradationer
--- blevet stillet saa eiendommeligt, at det næsten tager sig ud,
som om „Naturpocessen“ i „den nye Skabelse“ havde faaet% "Naturpocessen": sic, r dropped (printer's error). Quote balance: „staaer ... gjøres „en overnaturlig Begyndelse“ -- two openings, one closing, as printed (+1)
en betænkelig Overvægt over „Skabelsesprocessen“. Den Til\-%
stand, hvori Apostlene befandt sig, da de paa Pintsefesten
begyndte at tale med Tunger, betegnes nemlig som „en Til\-%
stand, hvori Bevidstheden ikke er sig selv mægtig, en Ekstasens,
en Henrykkelsens Tilstand, der var Udtrykket for de nye
aandelige Kræfters Gjennembrud og Rørelse i Sjælens Dyb,
og som mere havde Præget af en aandelig Naturtilstand, end
af det klare Bevidsthedsliv“; i Sandhed, en blandt den hellige
Histories mange „Skabelsesprocesser“ høist mærkværdig „Natur\-%
proces“! At denne „Ekstasens“ og „Henrykkelsens aandelige
Naturtilstand“ maa have varet temmelig længe og virket i
betydeligt Omfang, er især kjendeligt af den mærkværdige Ind\-%
flydelse, den har udøvet paa Forfatteren til Apostlenes Gjer\-%
ninger. Evangelisten Lukas er jo, utvivlsomt paa Grund af
en vis „aandelig Naturtilstand“, i sin Beretning kommen til
at stille det hele Mirakel paa Hovedet. Hvad Lukas fortæller
er, som bekjendt, dette, at de „gudfrygtige Mænd af alle
Folkeslag under Himmelen“ hørte Apostlene tale saaledes, at
de „forfærdedes og forundrede sig, og sagde til hverandre:
see, ere ikke alle disse, som tale, Galilæer? Og hvorledes
% --- p. 10 ---
\opage{10}\setcounter{footnote}{0}%
høre vi dem tale, hvert paa vort eget Maal, hvorudi vi ere
fødte? Parther og Meder og Elamiter, og vi, som boe i
Mesopothamia, Judæa og Cappadocia, Pontus og Asia, i
Phrygia og Pamphylia, Ægyptens og Lybiens Egne, ved
Cyrene, og vi her boende Romere, Jøder og Proselyter, Creter
og Araber? \emph{vi høre dem forkynde Guds store Gjer\-%
ninger med vore Tungemaal}“. Tydeligere, fuldstændigere
og høitideligere kunde Forfatteren ikke udtrykke sig, naar han
netop vilde lægge Vægt paa, at Underet just viste sig i
Sproggaven, og Sproggaven deri, at de eenfoldige Galilæer,
Mennesker, der ellers maatte antages kun at kunne eet eneste
Sprog, pludselig ved Aandens Indskydelse kom til at tale de
mange forskjellige fremmede Sprog. Dogmatiken, der har
udspeculeret Miraklet, veed imidlertid bedre Besked. „Den
christelige Dogmatik“ erklærer udtrykkelig, „at det ikke er
fornødent, her at tænke paa forskjellige Sprog, men paa den
Omstændighed, at „de“ (gudfrygtige Mænd af alle Folkeslag
under Himlen) „hver for sig følte sig tiltalte i deres egen
inderste \emph{Eiendommelighed}, at det indvortes Menneske i
Enhver saaledes tiltales af dette Henrykkelsens Tungemaal,
at Sjælen følte sig udløst af den vante, den naturlige Ind\-%
skrænkning og paa underfuld Maade blev sig selv aabenbar“.
Pintsemiraklet er da ved denne Forklaring bogstavelig blevet
endevendt. Det Mirakel, der i „Apostlenes Gjerninger“ skeer
med Apostlene, forklares i „Dogmatiken“ ved „en Tilstand,
hvori Bevidstheden ikke er sig selv mægtig, og bliver forsaavidt% quotation „en Tilstand ... opened here is not closed (+1); the „...aabenbar“ below closes without an opening (-1): both as printed
at ansee for et uegentligt Mirakel; medens de fremmede,
„gudfrygtige Mænds“ naturlige Glæde og Forundring over
hiint uegentlige Mirakel i „den christelige Dogmatik“ forvandles
til det egentlige Mirakel, det Mirakel nemlig, at Enhver af
de Tilstedeværende, endskjøndt han paa underfuld Maade blev
sig selv aabenbar“ --- saa at her altsaa i denne Henseende ikke
kunde være Tanke om Illusion --- ikke desmindre blev (natur\-%
ligviis ved Miraklet; thi Sligt er kun muligt ved et Mirakel)
sat i den Illusion, at han ved at høre et Tungemaal,
„hvorudi“ han vitterlig ikke var født, netop bildte sig ind at
høre „det Maal, hvorudi han var født“.

% --- p. 11 ---
\opage{11}\setcounter{footnote}{0}%
Dersom nu disse Prøver paa vor „christelige Dogma\-%
matiks“ dybsindige „Randsagelser“ og objective Forklaringer% "Dogmamatiks": sic, syllable "ma" doubled across the line break (printer's error)
gjør et lignende Indtryk paa Dem, Hr. Doctor, som paa
mig, saa blive vi snart enige om at opgive Haabet om den
formeentlig høiere Videnskab, den i Aanden forklarede Mirakel\-%
videnskab.

Det er Miraklet og atter Miraklet, der i strengeste For\-%
stand gjør Videnskaben Knuder. Vi kunne ikke løse disse
Knuder; lader os oprigtig tilstaae det: der gives ingen anden
Udvei! Dog, hvad siger jeg? en Udvei gives der jo dog, den
Udvei nemlig, at dække Knuderne til med et Par „værdifulde
Straaviske“ og overgyde „Straaviskene“ med Vrøvl. Dette
er unegtelig en Udvei, hvad Behandlingen angaaer. Men
der er det at sige herimod: Behandlingen er ikke videnskabelig;
og det Videnskabelige skulde jo ligge i Behandlingen.

Videnskaben er ikke religiøs, Religionen ikke viden\-%
skabelig; Astronomien er ikke andægtig, Andagten ikke
videnskabelig; Mathematiken er ikke bedende, Bønnen ikke
videnskabelig; Botaniken er uden Sakrament, Sakramentet er
ikke videnskabeligt. Men saa overbeviist jeg end er om, at
Tro og Viden, Religion og Videnskab maae --- til høiere
Forsoning --- begjære Skilsmisse, at altsaa Theologien, der
mæglende vil bære Kappen baade paa den religiøse og den
videnskabelige Skulder, maa bort; ligesaa overbeviist er jeg
om, at en Aabenbaring, der har „Livs-Ordet“ i sig selv,
ogsaa maa have Lyset i sig selv, at der, med andre Ord,
maa gives en religiøs Tænkning, grundig, klar, overbevisende,
en Tænkning, som ifølge sine eiendommelige Forudsætninger
hverken er eller kan eller vil være videnskabelig, men som dog
i Skarphed, Sikkerhed, Udholdenhed og Conseqvens veier op
imod Videnskaben.

De minder mig (Sendebr. S. 9) om Ørsteds Ord:
„Med Sandheden er det Gode uopløselig forbundet og den
\emph{fulde} Sandhed bringer selv sin Trøst med sig“, idet De med
Grund tilføier: „det er i Kraft af en Tro, en Tillid, altsaa
i Kraft af sin Subjectivitet, at han udtaler dette. Men
hvorledes skulde man komme til Erkjendelsen af den fulde
% --- p. 12 ---
\opage{12}\setcounter{footnote}{0}%
Sandhed, naar man ikke vilde \emph{respectere} det Factiske?
Saalænge noget i Erfaringen Givet bliver staaende og vidner
imod den Theorie, som man har dannet sig, da kan Theorien
ikke være sand. Men til det saaledes Givne kan ogsaa høre
Menneskehjertets \emph{Trang} til en levende Gud og til hans
Naade, saa at det kun i Troen derpaa kan finde Fred“.
Heri er jeg nu ganske og aldeles enig med Dem, Hr. Doctor,
ligesom vi jo ogsaa ere enige i, at „Gud, som \emph{Tankens
Product} ikke kan være Gjenstand for Tro“.

Men hvorledes skulle vi blive enige om det mærkværdige
Forhold, hvori De (Sendebr. S. 11) har stillet Troens Gud
til den menneskelige Fornuft? Det seer mig noget betænkeligt
ud. Skal jeg nemlig ved Fornuft og Fornuft, Tænkning og
Tænkning i begge Sætninger tænke det Samme, da er det
i Sandhed et vanskeligt Stykke. Lad mig gjentage, hvad
her staaer: Troens Gud kan, siger De, \emph{forsaavidt} (nemlig
som den Gud, der er Tankens Product, ikke kan være Gjen\-%
stand for Tro) være at kalde: „den sande Afgrund for den
menneskelige Fornuft“, og umiddelbart derpaa tilføier De:
„Men Troens Gud maa da (nemlig naar han er „den sande
Afgrund for den menneskelige Fornuft“) være Mere end dette,
Mere end den, ved hvem al Fornuft, al Tænkning maa
standse“... Hvorledes, spørger jeg, kan Troens Gud baade
være „den \emph{sande} Afgrund for den menneskelige Fornuft“ og
dog endnu Mere end den sande Afgrund, altsaa Afgrund i
den Grad, at han ikke er Afgrund; eller --- lad os nu rigtig
overveie, hvordan „det maa tænkes“ --- er Troens Gud maa\-%
skee en Afgrund, der standser den menneskelige Fornuft saa\-%
ledes, at Afgrunden selv standser for Fornuften, just som
Fornuften skulde standse for Afgrunden; eller er det maaskee
saaledes, at Afgrunden bliver Fornuft, idet Fornuften bliver
Afgrund; eller maaskee saaledes, at Fornuften først standser
for den sande Afgrund, naar den slet ikke standser; eller maa\-%
skee... Hr. Doctor, det Afgrundsstykke har en underlig Bi\-%
smag! det smager speculativt af „christelig Dogmatik“; jeg
kan ikke faae det ned.

% --- p. 13 ---
\opage{13}\setcounter{footnote}{0}%
De spørger (S. 12): Staaer Sandheden og Godheden% no opening „ is printed before "Staaer"; the quotation closes after "Gud?“ below (-1)
ikke virkeligt i et saa væsentligt Forhold til Gud, at man
slet ikke kan siges at have nogen praktisk Erkjendelse af Gud,
naar man ikke antager denne Forening af Sandhed og God\-%
hed med Gud?“ Jeg svarer: at den Troende ganske vist
maa antage, at Sandhed og Godhed ere forenede i Gud;
men at en Videnskab om denne Forening efter min Syns\-%
maade er umulig, sees ligefrem af den Modsætning mellem
Sandhedens og Godhedens Idee, hvorpaa min Opfattelse af
den principielle Forskjel imellem theoretisk og praktisk, viden\-%
skabelig og religiøs Erkjendelse er begrundet. Jeg overseer
ingenlunde Vexelvirkningen imellem Sandhed og Godhed,
mellem Tankens og Villiens Idealer. „Det er,“ har jeg jo
sagt, „en Selvfølge, at Tankens og Villiens Idealer maae
kunne indgaae i et ligesaa mangesidigt Vexelforhold som
Tanken og Villien selv; ethvert Forsøg paa at drage en
endelig bestemt Grædselinie, der skulde adskille den theoretiske% "Grædselinie": sic, n dropped (printer's error; cf. "Grændselinier", p. 3)
Erkjendelse fra den praktiske, er forsaavidt urimeligt og me\-%
ningsløst. Men derfor er her dog ligefuldt en Grændse, en
principiel Forskjel, der fremkommer derved, at alle Sjæle\-%
kræfter i Theorien underordnes Tænkningen, medens de i
Praxis underordnes Villien.“ (Om theoretisk og praktisk Er\-%
kjendelse S. 10). Min populære Fremstilling har De her
misforstaaet, som om jeg paa en udvortes Maade stillede
Villien op ved Siden af Tanken og med eensidig Opmærk\-%
somhed for det Praktiske kun havde Øie for Sætningen \textit{bonum
est, quia deus vult}, men glemte „at tilføie den anden
som ikke mindre berettiget: \textit{deus vult, quia bonum est}“.
Jeg vil i denne Anledning, Hr. Doctor, nødig bebyrde Dem
med metaphysiske Sentenser, thi jeg veed af Erfaring, at
Sligt ikke fører til Noget; men dersom De engang i et be\-%
leiligt Øieblik kom til at gjennemblade Grundideernes Logik
II. (S. 363--84), vilde De let overbevise Dem om, at jeg
over Sætningen: \textit{bonum est, quia deus vult}, ingenlunde
har glemt dens Modsætning: \textit{deus vult, quia bonum est}.

De mener, (Sendebr. S. 11) at jeg med Uret gjør
% --- p. 14 ---
\opage{14}\setcounter{footnote}{0}%
Religionen praktisk fra Først til Sidst, som om den ikke „med
noget væsentligt Baand var knyttet til en Viden og den
theoretiske Erkjendelse, som om i Religionen Tankerne slet
ikke betydede Noget i og for sig, men vare udelukkende prak\-%
tiske og Villien underordnede.“ Men naar De saa med denne
Forudsætning (Sendebr. S. 15) udtrykkelig henviser til
Skriften og erklærer, at „jo mere man stræber, søgende Næ\-%
ring for sin Tro, at leve sig ind i Skriften, desto mere op\-%
fordrer den os til \emph{Tænken} over de høieste Gjenstande, en
Tænken, ikke alene af praktisk, men ogsaa af theoretisk Art“;
da vedder jeg Ti imod Een, Hr. Doctor, at De her paa en
aldeles uklar Maade forvexler den \emph{religiøse Tænken}, om
De saa vil, den \emph{religiøse Theorie}, med den videnskabe\-%
lige. Hvor i al Verden er De dog kommen til at antage
den „Kundskab,“ den „Viden“ og den „Viisdom“, hvorom Chri\-%
stus og Apostlene tale, for videnskabelig Viden? Har Christus
da virkelig Charakteer af Videnskabsmand? Have maaskee
Isagogiken, Hermeneutiken, Exegetiken, Apologetiken og Dog\-%
matiken ved forenede Bestræbelser omsider bragt det derhen,
at Christi, af Dem selv citerede Ord (Matthæi 11., 25: „jeg
priser Dig Fader, at Du haver skjult dette for de Vise og
Forstandige og aabenbaret det for de Umyndige“), nu virkelig
lade sig opfatte som en Slags videnskabelig Indledning,
Indledning til Mirakelvidenskaben, den Videnskab, der god\-%
hedsfuld vil forklare os Miraklerne? Eller hvorledes har det
været Dem muligt i eet og samme Aandedræt at anføre de
to Steder af Paulus, først 1. Cor. 1, 14 „Herrens Ord hos
Propheten: Jeg vil forkaste de Vises Viisdom og tilintet\-%
gjøre de Forstandiges Forstand,“ og dernæst Bønnen for
Colossenserne, „at de maa fyldes med Guds Villies Erkjen\-%
delse i al Viisdom og aandelig Forstand“ (Col. 1, 9), og
her tænke paa videnskabelig Viisdom, videnskabelig Forstand?
Jeg seer nok, at De har udhævet Ordene: \emph{al Viisdom og
aandelig Forstand}; men jeg seer ogsaa, at De har for\-%
glemt at udhæve Ordene: \emph{Guds Villies Erkjendelse}.
Havde De lagt Mærke til, at den Erkjendelse, hvorom Paulus
% --- p. 15 ---
\opage{15}\setcounter{footnote}{0}%
taler, netop er \emph{Guds Villies} Erkjendelse, saa havde De
umulig kunnet oversee, at denne Villieserkjendelse trods al
Theori er og bliver af praktisk Art.

De siger (Sendebr. S. 17): „Foruden Kirkens og% a faint mark between "Kirkens" and "og" is raised spacing material, not a hyphen
Skriftens kjender man jo og \emph{Aandens}, den Hellig Aands
Vidnesbyrd. Skulde der ikke gives et \emph{videnskabeligt} Ud\-%
tryk for dette Vidnesbyrd, saa at man deri kunde have et
Princip vel overeensstemmende med Kirken og Skriften, men
ikke øst af disses Vidnesbyrd.“\footnote{Af Grundideernes Logik I. (Side 48--53; 239--43) og II. (Side
151--57; 343--47) vil det kunne skjønnes, at jeg fra Videns
Standpunkt ingenlunde har været uden Opmærksomhed for, hvad
Theologerne kalde \textit{testimonium spiritus}.} Denne Forventning vil% printed mark *) after the closing quote; the note stands at the foot of p. 15
blive skuffet. Det er jo dog en Kjendsgjerning, at den græsk-catholske% printed line break after the hyphen of "græsk-", a compound hyphen, kept
Theolog har beraabt sig paa den Hellig Aands
Vidnesbyrd ligesaavel som den romersk-catholske, den refor\-%
merte Theolog ligesaavel som den lutherske; og hvad har
Følgen været? At den Ene i Kraft af „Aandens Vidnes\-%
byrd“ har lyst Forbandelse over den Anden! Saaledes naar
vi see paa Kirkehistorien. See vi paa det Oprindelige, da
giver den Hellig Aand sig jo sit allerkraftigste Vidnesbyrd
ved Inspirationen. Hvilket er nu, ifølge Apostelen Peders
Tale (Apostl. Gjern. 2, 14--32) dette Vidnesbyrd? „Og det
skal skee i de sidste Dage, siger Gud, da vil jeg udgyde af
min Aand over alt Kjød; og Eders Sønner og Eders
Døttre skulle prophetere, og de Unge blandt Eder skulle see
Syner, og de Gamle skulle have Drømme“ o s. v.; men% "o s. v.": sic, no stop after "o"
et saadant Vidnesbyrd peger ikke i videnskabelig Retning.

Hvorledes den religiøse Erkjendelse fra Religionens
Synspunkt vil komme til at forholde sig til den videnskabe\-%
lige, er et stort og omfattende Problem, hvis Løsning maa
forsøges i en Religionsphilosophie; men en Religionsphilo\-%
sophie, der i nogen Grad skal løse Opgaven, maa da vide at
stille Videns Lys saaledes, at det, istedenfor at kaste forvirrende
Dæmringsskjær over Miraklerne, fra alle Sider kan belyse
% --- p. 16 ---
\opage{16}\setcounter{footnote}{0}%
Villiesprincipets og dermed Aabenbarings- og Troesprincipets
religiøse Selvstændighed.

Hvad jeg her i al Korthed har sammenskrevet til Svar
paa Deres ærede Sendebrev, beder jeg Dem nu, Hr. Doctor,
optage med Velvillie til nærmere Overveielse. Vi kunne, som
sagt, blive enige om meget væsentlige Punkter, dersom vi
blot kunde blive enige om, at en Videnskab, der forklarer
Mirakler, en Mirakelvidenskab, hører blandt de umulige
Ting.
% The letter ends here with a short centred rule; no signature, date or
% place is printed. PDF 19 verso and the rest of the scan are not text.
\begin{center}
\rule{0.25\textwidth}{0.4pt}
\end{center}

\end{document}
